% The fastless user manual. Build with "make doc", which runs pdflatex twice
% to produce doc/fastless.pdf. The reference manual page is doc/fastless.1;
% doc_tests.py checks that both describe every option.
\documentclass[11pt,a4paper]{article}
\usepackage[T1]{fontenc}
\usepackage[utf8]{inputenc}
\usepackage{textcomp}
\usepackage{charter}
\usepackage{tgheros}
\usepackage{lmodern}
\renewcommand{\rmdefault}{bch}
\renewcommand{\ttdefault}{lmtt}
\usepackage[a4paper,margin=25mm,top=28mm,bottom=28mm,headsep=10mm,headheight=14pt]{geometry}
\usepackage[protrusion=true,expansion=false]{microtype}
% Code keeps "--" as two hyphens.
\makeatletter
\expandafter\let\expandafter\fl@ttfamily\csname ttfamily \endcsname
\expandafter\def\csname ttfamily \endcsname{\fl@ttfamily\pdfnoligatures\font}
\makeatother
\usepackage{xcolor}
\usepackage{array}
\usepackage{booktabs}
\usepackage{longtable}
\usepackage{tabularx}
\usepackage{colortbl}
\usepackage{fancyhdr}
\usepackage{fancyvrb}
\usepackage{multicol}
\usepackage{parskip}
\usepackage{etoolbox}
\usepackage{caption}
\usepackage[hidelinks]{hyperref}

\definecolor{accent}{HTML}{1F5F8B}
\definecolor{accentlight}{HTML}{E8F0F6}
\definecolor{keybg}{HTML}{ECEEF1}
\definecolor{keyedge}{HTML}{B8BEC6}
\definecolor{rule}{HTML}{C9D3DC}
\definecolor{muted}{HTML}{5A6570}
\definecolor{good}{HTML}{2E7D4F}
\definecolor{warn}{HTML}{A15C00}

\hypersetup{
  colorlinks=true, linkcolor=accent, urlcolor=accent,
  pdftitle={fastless User Manual},
  pdfsubject={A less-compatible pager for very large and compressed files},
  pdfkeywords={pager, less, search, gzip, zstd, xz, bzip2},
  bookmarksnumbered=true, bookmarksopen=true,
}

% Headings: sans serif, in the accent color.
\makeatletter
\renewcommand\section{\@startsection{section}{1}{\z@}{-3.2ex \@plus -1ex \@minus -.2ex}%
  {1.6ex \@plus .2ex}{\sffamily\Large\bfseries\color{accent}}}
\renewcommand\subsection{\@startsection{subsection}{2}{\z@}{-2.6ex\@plus -1ex \@minus -.2ex}%
  {1ex \@plus .2ex}{\sffamily\large\bfseries\color{accent}}}
\renewcommand\subsubsection{\@startsection{subsubsection}{3}{\z@}{-2ex\@plus -1ex \@minus -.2ex}%
  {.6ex \@plus .2ex}{\sffamily\normalsize\bfseries}}
\renewcommand\@seccntformat[1]{\csname the#1\endcsname\hspace{.8em}}
\makeatother
\setcounter{tocdepth}{2}
% Long option names cannot be hyphenated; let their lines stretch instead.
\setlength{\emergencystretch}{3em}

\pagestyle{fancy}
\fancyhf{}
\fancyhead[L]{\sffamily\small\color{muted}fastless User Manual}
\fancyhead[R]{\sffamily\small\color{muted}\nouppercase{\leftmark}}
\fancyfoot[C]{\sffamily\small\thepage}
\renewcommand{\headrulewidth}{0.4pt}
\renewcommand{\headrule}{{\color{rule}\hrule height\headrulewidth width\headwidth}}
\renewcommand{\sectionmark}[1]{\markboth{#1}{}}

% Keys and other markup.
\newcommand{\kbd}[1]{{\setlength{\fboxsep}{1.2pt}\setlength{\fboxrule}{0.4pt}%
  \fcolorbox{keyedge}{keybg}{\vphantom{Ag}\small\ttfamily #1}}}
\newcommand{\cmd}[1]{\texttt{#1}}
% An option name may break after a hyphen inside it, as in --max-|mem, but
% not after its leading dashes.
\makeatletter
\newif\iffl@named
\def\fl@end{\fl@end}
\newcommand{\opt}[1]{\texttt{\fl@namedfalse\fl@walk#1\fl@end}}
\def\fl@walk#1{\ifx#1\fl@end\else\fl@char#1\expandafter\fl@walk\fi}
\def\fl@char#1{\ifx-#1-\iffl@named\discretionary{}{}{}\fi\else#1\fl@namedtrue\fi}
\makeatother
\newcommand{\file}[1]{\textit{#1}}
\newcommand{\env}[1]{\texttt{#1}}
\newcommand{\meta}[1]{\textit{#1}}
\newcommand{\fl}{\textsf{\small\color{accent}fastless}}
\newcommand{\bs}{\textbackslash}
\newcommand{\caret}{\textasciicircum}
\newcommand{\tl}{\raisebox{.42ex}{\texttildelow}}

% A shaded note.
\newsavebox{\notebox}
\newenvironment{note}[1][Note]%
  {\par\medskip\begin{lrbox}{\notebox}\begin{minipage}{\dimexpr\linewidth-14pt-3pt\relax}%
   \textsf{\bfseries\color{accent}#1.}\enspace}%
  {\end{minipage}\end{lrbox}\noindent{\setlength{\fboxsep}{7pt}\setlength{\fboxrule}{0pt}%
   \color{accent}\vrule width 3pt\colorbox{accentlight}{\color{black}\usebox{\notebox}}}\par\medskip}

% Shell examples.
\DefineVerbatimEnvironment{shell}{Verbatim}%
  {fontsize=\small,frame=leftline,framerule=2pt,rulecolor=\color{accent},
   xleftmargin=1em,framesep=1em,samepage=true,commandchars=\\\{\}}

% Tables of keys or options: a term column and a description.
\newenvironment{reftable}[1][0.30\linewidth]%
  {\renewcommand{\arraystretch}{1.25}%
   \begin{longtable}{@{}>{\raggedright\arraybackslash}p{#1}>{\raggedright\arraybackslash}p{\dimexpr\linewidth-#1-2\tabcolsep\relax}@{}}}%
  {\end{longtable}\vspace{-1ex}}
\newcommand{\thead}[2]{\toprule\textsf{\bfseries #1} & \textsf{\bfseries #2}\\\midrule\endhead\bottomrule\endfoot}
\newcommand{\cnt}{\textsuperscript{\color{accent}$\ast$}}

\begin{document}

% ---------------------------------------------------------------- title
\pagenumbering{alph}
\begin{titlepage}
\vspace*{40mm}
{\sffamily\fontsize{46}{50}\selectfont\bfseries\color{accent}fastless\par}
\vspace{6mm}
{\sffamily\LARGE A less-compatible pager for very large\\[1mm] and compressed files\par}
\vspace{10mm}
{\color{accent}\rule{\linewidth}{1.2pt}}
\vspace{6mm}

{\large User manual for version 1.0\par}
\vspace{2mm}
{\color{muted}October 2026\par}

\vfill
\begin{minipage}{0.78\linewidth}\small\color{muted}
\fl{} pages files of hundreds of gigabytes with the keys, options, and prompts of
GNU less. It maps files instead of reading them, decodes gzip, zstd, xz, and bzip2
on its own, and searches with every core.

\medskip
Copyright \textcopyright{} 2026 the fastless authors. fastless is free software,
distributed under the GNU General Public License, version 3.
\end{minipage}
\end{titlepage}

\pagenumbering{roman}
\tableofcontents
\clearpage
\pagenumbering{arabic}

% ---------------------------------------------------------------- intro
\section{Introduction}

\fl{} is a terminal pager in the tradition of \cmd{more} and GNU \cmd{less}. If
you know \cmd{less}, you already know how to use it: the keys, the command-line
options, the \env{LESS} variable, prompts, marks, lesskey files, and input
preprocessors all work the same way. What is different is the scale it is built
for.

\begin{itemize}
\item \textbf{Huge files open instantly.} A plain file is mapped into memory
  read-only. Nothing is copied and no line index is built, so a 200\,GB log opens
  as fast as a small one, and memory use does not grow with the file.
\item \textbf{Searches use every core.} A search that runs further than 8\,MiB is
  split into pieces that are searched in parallel; the first match in file order
  is still the one reported. Literal searches use AVX2 instructions on CPUs that
  have them.
\item \textbf{Compressed files are understood natively.} gzip, zstd, xz, and
  bzip2 are recognised by their content, decoded (on several threads for xz and
  bzip2), and verified before the first page is shown.
\item \textbf{Three ways to search.} Besides regular expressions, \fl{} can
  search for literal bytes, which is the fastest, or match whole lines against a
  glob pattern.
\item \textbf{Pipes are paged as they arrive}, just as in \cmd{less}, without
  slowing down ordinary files.
\end{itemize}

Section~\ref{sec:less} lists everything that behaves differently from \cmd{less}.

% ---------------------------------------------------------------- start
\section{Getting started}

\subsection{Building and installing}

\fl{} is written in C and builds with GCC or Clang on 64-bit Linux and macOS. It
needs \cmd{pkg-config}, libzstd, liblzma, and a gzip decoder. Clone the
repository with its submodules, which supply zlib-ng and the lbzip2 decoder:

\begin{shell}
git clone --recursive https://github.com/dvquilting/fastless.git
cd fastless
make
sudo make install          \textit{# into /usr/local; PREFIX=... to change}
\end{shell}

On Debian and Ubuntu, the build dependencies are:

\begin{shell}
sudo apt install build-essential pkg-config cmake libisal-dev zlib1g-dev \bs
                 libbz2-dev libzstd-dev liblzma-dev
\end{shell}

For gzip, the build uses Intel ISA-L when it is installed, otherwise the bundled
zlib-ng, otherwise the system zlib; \cmd{make ZLIB=isal}, \cmd{ZLIB=ng}, or
\cmd{ZLIB=system} chooses one. bzip2 uses the bundled lbzip2 decoder, or the
system libbz2 with \cmd{make BZIP2=system}. \cmd{make check} runs the test
suites, and \cmd{make doc} rebuilds this manual.

\subsection{A first look}

\begin{shell}
fastless huge.log                  \textit{# page a file}
fastless huge.log.zst other.log.gz \textit{# compressed files, two of them}
some-command | fastless            \textit{# a pipe, shown as it arrives}
LESS=-iR fastless build.log        \textit{# options from the environment}
fastless +G --literal app.log      \textit{# start at the end, literal searches}
\end{shell}

Inside the pager, \kbd{Space} and \kbd{b} move a window forward and back,
\kbd{j} and \kbd{k} a line, \kbd{g} and \kbd{G} go to the start and the end, and
\kbd{/} searches. \kbd{h} shows a summary of every command, and \kbd{q} quits.

The status line at the bottom is \fl's own. It shows the file name, the byte
position and percentage, the search mode, and a reminder of the main keys:

\begin{shell}
huge.log | byte 1048576/214748364800 (0.0%) | regex | q quit; h help; / search
\end{shell}

\cmd{-m}, \cmd{-M}, and \cmd{-P} switch to \cmd{less}'s prompts instead
(Section~\ref{sec:prompts}).

When standard output is not a terminal, \fl{} does not page at all: it copies its
input to standard output, decoded if it was compressed. So
\cmd{fastless huge.log.zst | grep ERROR} works like \cmd{zstdcat}.

% ---------------------------------------------------------------- moving
\section{Moving around}

Commands are single keys, typed without \kbd{Return}. Those marked \cnt{} may
be preceded by a number $N$: \cmd{10j} moves ten lines, \cmd{50\%} goes halfway,
and \cmd{3n} finds the third match. \kbd{\caret X} means Control-X, and
\kbd{ESC-x} means Escape followed by x.

\subsection{Scrolling}

\begin{reftable}
\thead{Keys}{Action}
\kbd{Space} \kbd{f} \kbd{\caret F} \kbd{\caret V} \kbd{PgDn}\cnt & Forward one window, or $N$ lines. \\
\kbd{b} \kbd{\caret B} \kbd{ESC-v} \kbd{PgUp}\cnt & Backward one window, or $N$ lines. \\
\kbd{z} \,/\, \kbd{w}\cnt & Forward / backward a window, setting the window size to $N$. \\
\kbd{ESC-Space}\cnt & Forward a window, even at the end of the file. \\
\kbd{j} \kbd{e} \kbd{Return} \kbd{\caret E} \kbd{\caret N} \kbd{Down}\cnt & Forward one line, or $N$ lines. \\
\kbd{k} \kbd{y} \kbd{\caret Y} \kbd{\caret P} \kbd{\caret K} \kbd{Up}\cnt & Backward one line, or $N$ lines. \\
\kbd{J} \,/\, \kbd{K} \kbd{Y}\cnt & Forward / backward a line, even past either end. \\
\kbd{d} \kbd{\caret D} \,/\, \kbd{u} \kbd{\caret U}\cnt & Forward / backward half a window, setting the half-window to $N$. \\
\kbd{ESC-)} \kbd{Right} \,/\, \kbd{ESC-(} \kbd{Left}\cnt & Right / left half a screen, or $N$ columns. \\
\kbd{ESC-\}} \kbd{\caret Right} & Right to the end of the longest line on the screen. \\
\kbd{ESC-\{} \kbd{\caret Left} & Back to the first column. \\
\kbd{F} & Follow the end of the file, like \cmd{tail -f}. Any key stops. \\
\kbd{ESC-F} & Follow, stopping when the search pattern appears. \\
\kbd{r} \kbd{\caret R} \kbd{\caret L} & Repaint the screen. \\
\kbd{R} & Repaint, first taking in growth of the file. \\
\end{reftable}

Long lines are never wrapped; they end with a \texttt{>} in the last column, and
the horizontal commands bring the rest into view.

\subsection{Jumping}

\begin{reftable}
\thead{Keys}{Action}
\kbd{g} \kbd{<} \kbd{ESC-<} \kbd{Home}\cnt & The first line, or line $N$. \\
\kbd{G} \kbd{>} \kbd{ESC->} \kbd{End}\cnt & The last line, or line $N$. \\
\kbd{p} \kbd{\%}\cnt & $N$ percent into the file; decimals such as \cmd{37.5\%} work. \\
\kbd{P}\cnt & Byte offset $N$. \hfill\textsf{\footnotesize\color{accent}fastless} \\
\kbd{\{} \kbd{(} \kbd{[}\cnt & The bracket closing the first such bracket on the top line. \\
\kbd{\}} \kbd{)} \kbd{]}\cnt & The bracket opening the last such bracket on the bottom line. \\
\kbd{ESC-\caret F}\,\meta{a b} & The same for any pair of characters \meta{a} and \meta{b}; \kbd{ESC-\caret B} backward. \\
\end{reftable}

\subsection{Marks}

\begin{reftable}
\thead{Keys}{Action}
\kbd{m}\meta{x} \,/\, \kbd{M}\meta{x} & Mark the top / bottom line with the letter \meta{x}. \\
\kbd{'}\meta{x} \,or\, \kbd{\caret X\caret X}\meta{x} & Go to mark \meta{x}. \\
\kbd{'}\kbd{'} & Back to the previous position, before the last jump. \\
\kbd{'}\kbd{\caret} \,/\, \kbd{'}\kbd{\$} & The start / end of the file. \\
\kbd{ESC-m}\meta{x} & Clear mark \meta{x}. \\
\end{reftable}

With \opt{--save-marks}, marks are kept in the history file between sessions.
With \opt{--mouse}, the wheel scrolls (\opt{--wheel-lines} at a time) and a click
sets mark \cmd{\#}.

% ---------------------------------------------------------------- search
\section{Searching}

\subsection{Search commands}

\begin{reftable}
\thead{Keys}{Action}
\kbd{/}\meta{pattern}\cnt & Search forward for the ($N$-th) line with a match. \\
\kbd{?}\meta{pattern}\cnt & Search backward. \\
\kbd{n} \,/\, \kbd{N}\cnt & Repeat the search in the same / opposite direction. \\
\kbd{ESC-n} \,/\, \kbd{ESC-N}\cnt & Repeat, continuing into the following / preceding files. \\
\kbd{\&}\meta{pattern} & Show only the lines that match; \kbd{\&} alone shows all lines again. \\
\kbd{ESC-u} \,/\, \kbd{ESC-U} & Toggle highlighting / forget the pattern. \\
\kbd{ESC-r} \kbd{ESC-l} \kbd{ESC-w} & Regular-expression, literal, or glob mode. \hfill\textsf{\footnotesize\color{accent}fastless} \\
\end{reftable}

A match is reported by its byte offset, and the line holding it is placed on the
jump target line (\opt{-j}; the top line by default). Every match on the screen is
highlighted, overlapping ones included; \opt{-g} highlights only the one found and
\opt{-G} none. A new forward search starts at the target line; \opt{-a} starts it
after the screen, and \opt{-A} just after the target line. \opt{--incsearch}
searches as you type, each keystroke cancelling the search still running.

\subsection{Three search modes}

\begin{reftable}[0.22\linewidth]
\thead{Mode}{Matches}
\textbf{Regex} \newline \opt{--regex}, \kbd{ESC-r} \newline \textsf{\footnotesize\color{muted}the default} &
POSIX extended regular expressions, such as \cmd{\caret ERROR.*timeout}, matched
within each line without its newline. \\
\textbf{Literal} \newline \opt{--literal}, \kbd{ESC-l} &
The pattern's exact bytes, anywhere, even across line ends. The fastest mode,
using AVX2 where available. \\
\textbf{Glob} \newline \opt{--glob}, \kbd{ESC-w} &
A pattern for the \emph{whole line}: \cmd{*} matches any bytes, \cmd{?} one byte,
and \cmd{\bs} quotes the next character. Use \cmd{*text*} to find \meta{text}
anywhere in a line. Bracket classes are not supported. \\
\end{reftable}

Switching modes recompiles the current pattern, so \kbd{n} goes on with the same
text under the new rules. A regex or glob is not run on every line: \fl{} first
works out the longest run of plain characters that every match must contain,
finds the lines holding it with the literal matcher, and runs the full pattern
only on those. For \cmd{\caret ERROR.*timeout}, only lines containing
\cmd{ERROR} are examined. For alternatives such as \cmd{timeout|refused}, a run
from each alternative is looked for, and only lines holding one are examined.

\begin{note}[Tip]
For the fastest searches of very large files, use literal mode, and give a
starting point with a mark, a percentage, or \kbd{P} so that less of the file
has to be read from storage.
\end{note}

\subsection{Case}

\opt{-i} ignores case when the pattern contains no capital letters (``smart
case''), \opt{-I} always ignores it, and \opt{--case-sensitive} never does.
Matching is done in bytes in the C locale: case folding covers ASCII letters
only, and a UTF-8 pattern matches its encoded bytes.

\subsection{Pattern modifiers}

As in \cmd{less}, a pattern may begin with modifiers, which also change the
prompt's label:

\begin{reftable}[0.18\linewidth]
\thead{Modifier}{Effect}
\kbd{\caret N} or \cmd{!} & Find lines that do \emph{not} match. \\
\kbd{\caret E} or \cmd{*} & Continue through the following (or preceding) files. \\
\kbd{\caret F} or \cmd{@} & Start at the first file (\cmd{/}) or the last (\cmd{?}). \\
\kbd{\caret K} & Highlight the matches without moving. \\
\kbd{\caret R} & Search literally, for this pattern only. \\
\kbd{\caret S}\,\meta{n} & Match parenthesised subexpression \meta{n} and highlight just it. \\
\kbd{\caret W} & Wrap to the other end of the file if nothing is found. \\
\kbd{\caret L} & Enter the next character literally. \\
\end{reftable}

The punctuation forms \cmd{!}, \cmd{*}, and \cmd{@} are modifiers only in regex
mode, because glob patterns such as \cmd{*error*} and literal ones such as
\cmd{!important} begin with them. \opt{--search-options=EFKNRW} sets modifiers for
every search. \kbd{n} and \kbd{N} keep the modifiers except \kbd{\caret F} and
\kbd{\caret K}, and an empty pattern reuses the current one with new modifiers.

\subsection{Filtering}

\kbd{\&}\meta{pattern} shows only matching lines (with \kbd{\caret N}, only
those that do not match), and movement, searches, and marks then refer to the
lines shown. Filtering keeps no index: each line shown is found with the same
threaded search, so a page of sparse matches costs a scan of the bytes between
them.

\subsection{Searching without the pager}

\opt{--find} prints the zero-based byte offset of the first match and exits, and
with \opt{--all}, every match, overlapping ones included. \opt{--offset} starts
the search at a byte position. Glob mode prints one offset per matching line.

\begin{shell}
fastless --find '\caret ERROR.*timeout' huge.log
fastless --literal --find needle --all huge.log.zst
fastless --literal --find needle --offset 150000000000 huge.log
\end{shell}

The exit status is 0 if something matched, 1 if nothing did, 2 for an error,
and 130 if the search was interrupted.

% ---------------------------------------------------------------- input
\section{Files, compressed input, and pipes}

\subsection{Several files}

\begin{reftable}
\thead{Keys}{Action}
\kbd{:e} [\meta{file}] \kbd{\caret X\caret V} & Examine a new file; \cmd{:e} alone rereads the current one. \\
\kbd{:n} \,/\, \kbd{:p}\cnt & The next / previous file named on the command line. \\
\kbd{:x}\cnt & The first (or $N$-th) file. \\
\kbd{:d} & Remove the current file from the list. \\
\kbd{=} \kbd{\caret G} \kbd{:f} & Show the file name, position, and size. \\
\end{reftable}

Only the current file is kept open. Moving to another file releases a decoded
file's memory, and coming back decodes it again. Standard input, pipes, and
devices cannot be read twice, so they stay open with what was read.

\subsection{Compressed files}

gzip, zstd, xz, and bzip2 input is recognised by its first bytes, so names and
extensions do not matter, and \opt{--compression} can force a format or
\cmd{none}. Concatenated gzip members, zstd frames, xz streams, and bzip2 streams
are all read; only zero bytes may follow the last one. The whole file is decoded
and its checksums verified before the first page appears. Progress is shown on
standard error, and \kbd{\caret C} cancels. Damaged or truncated data is an
error, never a partial file.

\begin{reftable}[0.15\linewidth]
\thead{Format}{Decoder}
gzip & Intel ISA-L or zlib-ng; one thread. \\
zstd & libzstd; one thread, windows up to 128\,MiB. \\
xz & liblzma; several threads when the file has several blocks, as threaded encoders write. \\
bzip2 & lbzip2's block decoder on every thread, or libbz2 on one. \\
\end{reftable}

Decoded data is kept in memory up to \opt{--max-mem} (10\,GiB unless set; fractions
such as \cmd{0.5} work). Larger output moves to a temporary file, created with mode
0600 in \opt{--tmpdir}, \env{\$TMPDIR}, or \file{/tmp} and deleted at once, so it
disappears when \fl{} exits for any reason.

\begin{note}[Plan for the disk space]
Decoded data can be many times the size of the compressed file, and the temporary
file needs room for all of it. On Linux, \file{/tmp} is often held in memory:
for large files, point \opt{--tmpdir} at a disk.
\end{note}

\subsection{Pipes and standard input}

Without a file name, or with \cmd{-}, \fl{} reads standard input and takes keys
from \file{/dev/tty}. Named pipes open like files; devices need \opt{-f}.

Streams are shown while they are still being written, so
\cmd{make 2>\&1 | fastless} and \cmd{tail -f log | fastless} show output as it
arrives. While a stream is being read, its size is not known: the status shows
\textit{(reading)}, moving past the end waits for more data, \kbd{G} and
\kbd{\%} follow the end until the input ends, and forward searches go on into
input as it arrives. Any key ends a wait.

Up to \opt{--max-mem} (but at least 64\,MiB) is read ahead so that the writer is
never held up. Beyond that, input is read only as far as the screen, a search,
or a jump needs, so \cmd{yes | fastless} does not fill the disk. A compressed
stream is read whole, then decoded. When standard output is not a terminal, a
pipe is copied straight through in constant memory.

\subsection{Input preprocessors}

\env{LESSOPEN} and \env{LESSCLOSE} work as in \cmd{less}: \cmd{|cmd~\%s} pages the
command's output, \cmd{||cmd~\%s} also accepts empty output, and \cmd{cmd~\%s}
names a replacement file. Files that begin with a compressed format's magic bytes
skip the preprocessor, because the built-in decoders are faster;
\cmd{-L --compression=none} shows their raw bytes. \opt{-L} ignores
\env{LESSOPEN} altogether, and \opt{--show-preproc-errors} reports failures.

% ---------------------------------------------------------------- display
\section{Display, line numbers, and prompts}\label{sec:prompts}

\subsection{How text is shown}

Text is shown as UTF-8 when \env{LESSCHARSET} or the locale says UTF-8;
otherwise bytes above 127 appear as \cmd{<XX>}. Control characters appear in bold
as \cmd{\caret X}, and invalid or unprintable characters as \cmd{<XX>} or
\cmd{<U+XXXX>}, so a file can never send escape sequences to the terminal.
\opt{-R} passes ANSI colour sequences through. As in \cmd{less}, an overstruck
character is bold, an underscore overstrike underlines, a carriage return before
a newline is hidden, and tabs expand to the \opt{-x} stops.

\subsection{Line numbers}

\opt{-N} shows a line number before each line. Prompts can show line numbers too:
\opt{-M} and \kbd{=} give \textit{lines 6--12/100}. Nothing is counted unless
something on the screen shows a number, so without \opt{-N} or such a prompt
there is no cost at all. When numbers are needed, newlines are counted with SIMD
on all search threads and remembered per 64\,MiB piece. A long jump into
uncounted territory never waits: the screen appears at once, the numbers fill in
when they are known, and a counter shows the progress of longer counts.
\kbd{\caret C} during counting turns line numbers off, as in \cmd{less}.

\subsection{Prompts}

\opt{-m}, \opt{-M}, and \opt{-P} select \cmd{less}'s medium, long, or custom
prompts and its prompt language. \opt{-Ps} (or \opt{-P} alone) sets the short
prompt, \opt{-Pm} and \opt{-PM} the medium and long ones, \opt{-P=} the \kbd{=}
message, \opt{-Ph} the help prompt, and \opt{-Pw} the prompt while waiting for a
pipe. The escapes are:

\begin{multicols}{2}\small
\begin{tabular}{@{}>{\ttfamily}l l@{}}
\%f & file name \\
\%F & last component of the name \\
\%g & name, quoted for the shell \\
\%i, \%m & file index, number of files \\
\%x & next file's name \\
\%b, \%B & byte offset, size \\
\%p, \%P & percent, by bytes / by lines \\
\end{tabular}
\columnbreak

\begin{tabular}{@{}>{\ttfamily}l l@{}}
\%l, \%L & line number, last line \\
\%d, \%D & page number, page count \\
\%c & column of the left edge \\
\%s & size, as \%B \\
\%t, \%T & trim spaces / the word ``file'' \\
\%E & the editor \\
?\textit{x}\ldots:\ldots. & conditionals \\
\end{tabular}
\end{multicols}

\vspace{-1ex}
A line-number item shows \cmd{?}, and its conditional is false, until the number
is known. While a pipe is still being read, the conditions on its size
(\cmd{?s}, \cmd{?B}, \cmd{?p}, \cmd{?P}, \cmd{?L}, \cmd{?D}, \cmd{?e}) are false.

% ---------------------------------------------------------------- misc commands
\section{Other commands}

\begin{reftable}
\thead{Keys}{Action}
\kbd{-}\meta{x} \,/\, \kbd{-}\kbd{-}\meta{name} & Toggle or set an option; \cmd{-+}\meta{x} resets it, \cmd{-!}\meta{x} sets its opposite. \\
\kbd{\_}\meta{x} \,/\, \kbd{\_}\kbd{\_}\meta{name} & Show an option's setting. \\
\kbd{+}\meta{cmd} & Run \meta{cmd} in each file examined from now on. \\
\kbd{!}\meta{command} & Run a shell command: \cmd{\%} is the file name, \cmd{\#} the previous file, \cmd{!!} the last command. \\
\kbd{\#}\meta{command} & The same, expanding prompt escapes. \\
\kbd{|}\meta{x}\,\meta{command} & Pipe the text from mark \meta{x} (or \cmd{\caret}, \cmd{\$}, \cmd{.}) through the screen into \meta{command}. \\
\kbd{s}\,\meta{file} & Save the input; for compressed files, the decoded data. \\
\kbd{v} & Edit the file with \env{\$VISUAL}, \env{\$EDITOR}, or \cmd{vi}, at the current line. \\
\kbd{h} \kbd{H} \,/\, \kbd{V} & Help / version. \\
\kbd{q} \kbd{Q} \kbd{:q} \kbd{:Q} \kbd{ZZ} & Quit, or leave the help. \\
\kbd{\caret C} & Cancel a search, count, wait, or \kbd{F}. \\
\kbd{\caret Z} & Suspend. \\
\end{reftable}

\subsection{Editing a prompt}

Prompts for patterns, file names, and commands have \cmd{less}'s line editing:
\kbd{Left} \kbd{Right} (\kbd{ESC-h} \kbd{ESC-l}), word moves with \kbd{\caret Left}
\kbd{\caret Right} (\kbd{ESC-b} \kbd{ESC-w}), \kbd{Home} \kbd{End} (\kbd{ESC-0}
\kbd{ESC-\$}), \kbd{Del} (\kbd{ESC-x}), word deletion with \kbd{ESC-Bksp} and
\kbd{ESC-X}, \kbd{\caret U} to clear the line, \kbd{\caret V} to enter a character
literally, and \kbd{Up} \kbd{Down} for history, matching what was typed.
\kbd{Bksp} on an empty prompt, \kbd{\caret G}, or a lone \kbd{ESC} cancels.
History is kept in \file{\tl/.fastlesshst} (\env{FASTLESSHISTFILE}).

% ---------------------------------------------------------------- options
\section{Options}

Options are read from the \env{LESS} environment variable first, then from the
command line. Long options may be shortened to any unique prefix, and most can be
changed inside the pager with \kbd{-}. In \env{LESS}, string values end at
\cmd{\$} and unknown options are ignored. Options marked
\textsf{\footnotesize\color{accent}fastless} are not in \cmd{less}.

\newcommand{\own}{\hfill\textsf{\footnotesize\color{accent}fastless}}

\subsection{Searching}
\begin{reftable}[0.34\linewidth]
\thead{Option}{Meaning}
\opt{--regex}, \opt{--literal}, \opt{--glob} & The search mode. \own \\
\opt{-i}, \opt{--ignore-case} & Ignore case if the pattern has no capitals. \\
\opt{-I}, \opt{--IGNORE-CASE} & Always ignore case. \\
\opt{--case-sensitive} & Turn \opt{-i} and \opt{-I} off. \own \\
\opt{-a}, \opt{--search-skip-screen} & Start forward searches after the screen. \\
\opt{-A}, \opt{--SEARCH-SKIP-SCREEN} & Start them just after the target line. \\
\opt{-g}, \opt{--hilite-search} & Highlight only the match found. \\
\opt{-G}, \opt{--HILITE-SEARCH} & Highlight nothing. \\
\opt{--incsearch} & Search while the pattern is typed. \\
\opt{--search-options=}\meta{EFKNRW} & Modifiers for every search; \cmd{-} clears those before it. \\
\opt{-p}\,\meta{pattern}, \opt{--pattern} & Search when the first file opens. \\
\opt{-j}\,\meta{n}, \opt{--jump-target} & The target line: a line (negative from the bottom) or a fraction such as \cmd{.5}. \\
\opt{--match-shift=}\meta{n} & Columns kept left of a match scrolled into view. \\
\opt{--find=}\meta{pattern} & Print the first match's byte offset and exit. \own \\
\opt{--all} & With \opt{--find}, print every match. \own \\
\opt{--offset=}\meta{byte} & Start the pager or the search at a byte offset. \own \\
\end{reftable}

\subsection{Display}
\begin{reftable}[0.34\linewidth]
\thead{Option}{Meaning}
\opt{-R}, \opt{--RAW-CONTROL-CHARS} & Pass ANSI colour sequences through. \\
\opt{-r}, \opt{--raw-control-chars} & The same as \opt{-R}. \\
\opt{-N}, \opt{--LINE-NUMBERS} & Show line numbers. \\
\opt{-n}, \opt{--line-numbers} & Show no line numbers, in prompts either. \\
\opt{--line-num-width=}\meta{n} & Width of the number column (7). \\
\opt{--no-number-headers} & Leave header lines unnumbered. \\
\opt{-s}, \opt{--squeeze-blank-lines} & Show runs of blank lines as one. \\
\opt{-x}\,\meta{n,...}, \opt{--tabs} & Tab stops (8). \\
\opt{-\tl}, \opt{--tilde} & No \cmd{\tl} after the end of the file. \\
\opt{-\#}\,\meta{n}, \opt{--shift} & Horizontal step: columns, a fraction, or 0 for half the screen. \\
\opt{--header=}\meta{L,C,N} & Keep \meta{L} lines (from line \meta{N}) and \meta{C} columns in place. \\
\opt{--no-search-header-lines} & Searches skip header lines; \opt{--no-search-header-columns} columns, \opt{--no-search-headers} both. \\
\opt{-J}, \opt{--status-column} & A column marking marks and matching lines. \\
\opt{--status-col-width=}\meta{n} & Its width (2). \\
\opt{--status-line} & Highlight marked lines. \\
\opt{-u}, \opt{--underline-special} & Show backspaces and carriage returns as control characters. \\
\opt{-U}, \opt{--UNDERLINE-SPECIAL} & Also tabs. \\
\opt{--proc-backspace}, \opt{--proc-return}, \opt{--proc-tab} & Process that character whatever \opt{-u} says. \\
\opt{--PROC-BACKSPACE}, \opt{--PROC-RETURN}, \opt{--PROC-TAB} & Show it as a control character whatever \opt{-u} says. \\
\opt{--rscroll=}\meta{c} & The mark for a chopped line (\cmd{>}; \cmd{-} for none). \\
\opt{-w}, \opt{--hilite-unread} & Highlight the first new line after a page forward. \\
\opt{-W}, \opt{--HILITE-UNREAD} & \ldots{} after any forward movement. \\
\opt{-m}, \opt{--long-prompt} & \cmd{less}'s medium prompt. \\
\opt{-M}, \opt{--LONG-PROMPT} & \cmd{less}'s long prompt. \\
\opt{-P}\meta{prompt}, \opt{--prompt} & A custom prompt (Section~\ref{sec:prompts}). \\
\opt{-c}, \opt{--clear-screen}, \opt{-C}, \opt{--CLEAR-SCREEN} & Repaint by clearing the screen. \\
\opt{-X}, \opt{--no-init} & Use the normal screen, not the alternate one. \\
\opt{--redraw-on-quit} & Leave the last screen on the terminal. \\
\opt{-z}\,\meta{n}, \opt{--window} & Window size; negative counts from the screen height. \\
\opt{--mouse}, \opt{--MOUSE} & Wheel scrolling; upper case reverses it. \\
\opt{--wheel-lines=}\meta{n} & Lines per wheel step (1). \\
\end{reftable}

\subsection{Quitting and the bell}
\begin{reftable}[0.34\linewidth]
\thead{Option}{Meaning}
\opt{-e}, \opt{--quit-at-eof} & Quit on a forward command at the end. \\
\opt{-E}, \opt{--QUIT-AT-EOF} & Quit as soon as the end is shown. \\
\opt{-F}, \opt{--quit-if-one-screen} & Print a file that fits on one screen and exit. \\
\opt{-K}, \opt{--quit-on-intr} & Quit with status 2 on \kbd{\caret C}. \\
\opt{--intr=}\meta{c} & Another character that acts as \kbd{\caret C} (\kbd{\caret X}). \\
\opt{-q}, \opt{--quiet}, \opt{--silent} & Ring the bell only for errors. \\
\opt{-Q}, \opt{--QUIET}, \opt{--SILENT} & Never ring the bell. \\
\end{reftable}

\subsection{Input and resources}
\begin{reftable}[0.34\linewidth]
\thead{Option}{Meaning}
\opt{--compression=}\meta{format} & \cmd{auto}, \cmd{none}, \cmd{gzip}, \cmd{zstd}, \cmd{xz}, or \cmd{bzip2}. \own \\
\opt{--max-mem=}\meta{GiB} & Memory for decoded or piped input before a temporary file (10). \own \\
\opt{--tmpdir=}\meta{dir} & Where temporary files go. \own \\
\opt{--threads=}\meta{n} & Threads for searching, counting, xz and bzip2: 1--256, or 0 for one per CPU up to 32. \own \\
\opt{--simd=}\meta{mode} & \cmd{auto}, \cmd{off}, or \cmd{avx2} for literal searches. \own \\
\opt{-f}, \opt{--force} & Open devices and other non-regular files. \\
\opt{-L}, \opt{--no-lessopen} & Ignore \env{LESSOPEN}. \\
\opt{--show-preproc-errors} & Report a failing preprocessor. \\
\opt{--follow-name} & \kbd{F} reopens a file replaced under its name. \\
\opt{--lesskey-src=}\meta{file} & Key bindings from a lesskey source file. \\
\opt{--lesskey-content=}\meta{text} & Key bindings as text; lines end at newlines or \cmd{;}. \\
\opt{--save-marks} & Keep marks between sessions. \\
\opt{--no-histdups} & Drop earlier copies of a command from the history. \\
\opt{+}\meta{cmd}, \opt{++}\meta{cmd} & Run \meta{cmd} in the first / every file, as in \cmd{+G} or \cmd{+/error}. \\
\opt{-V}, \opt{--version} & Print the version. \\
\opt{-?}, \opt{--help} & Print an option summary. \\
\end{reftable}

\subsection{Accepted, but without effect}

\opt{-b} (\opt{--buffers}), \opt{-B} (\opt{--auto-buffers}), \opt{-d}
(\opt{--dumb}), \opt{-h} (\opt{--max-back-scroll}), \opt{-S}
(\opt{--chop-long-lines}), \opt{-y} (\opt{--max-forw-scroll}), \opt{--file-size},
\opt{--no-keypad}, \opt{--no-vbell}, and \opt{--exit-follow-on-close} are accepted
so that existing \env{LESS} settings keep working. Lines are always chopped, as
with \opt{-S}.

\subsection{Not supported}

\opt{-o}, \opt{--log-file}, \opt{-O}, and \opt{--LOG-FILE} (log files);
\opt{-t}, \opt{--tag}, \opt{-T}, and \opt{--tag-file} (tags); \opt{-D},
\opt{--color}, and \opt{--use-color} (colours); \opt{--wordwrap};
\opt{--modelines}; \opt{-"} and \opt{--quotes}; \opt{--use-backslash}; and
\opt{-k} and \opt{--lesskey-file} (compiled lesskey files) are refused with a
message.

% ---------------------------------------------------------------- keys / env
\section{Key bindings and the environment}

\subsection{lesskey files}

Keys can be rebound with \cmd{lesskey} source: \opt{--lesskey-src}, or else the
first of \env{\$LESSKEYIN}, \file{\$XDG\_CONFIG\_HOME/lesskey},
\file{\tl/.config/lesskey}, and \file{\tl/.lesskey}; then \env{LESSKEY\_CONTENT}
and \opt{--lesskey-content}. The \cmd{\#command}, \cmd{\#line-edit}, and
\cmd{\#env} sections, \cmd{\#stop}, \cmd{\bs k} key names, and extra strings
work as in \cmd{lesskey}(1). Besides \cmd{less}'s actions there are
\cmd{regex-mode}, \cmd{literal-mode}, and \cmd{glob-mode}; the tag and OSC~8
hyperlink actions are accepted but do nothing.

\begin{shell}
#command
t     forw-screen
x     literal-mode
#env
LESS = -iR
\end{shell}

\subsection{Environment variables}

\begin{reftable}[0.30\linewidth]
\thead{Variable}{Use}
\env{LESS} & Options read before the command line. \\
\env{LESSOPEN}, \env{LESSCLOSE} & Input preprocessors. \\
\env{LESSEDIT} & The editor command for \kbd{v}, in prompt language (\cmd{\%E ?lm+\%lm. \%g}). \\
\env{VISUAL}, \env{EDITOR} & The editor; empty counts as unset (\cmd{vi}). \\
\env{SHELL} & The shell for \kbd{!}, \kbd{\#}, and \kbd{|}. \\
\env{LESSSECURE} & When \cmd{1}: no shell commands, pipes, saving, editing, \cmd{:e}, preprocessors, or history. \\
\env{LESSCHARSET} & \cmd{utf-8} for UTF-8, anything else for single bytes; unset, the locale decides. \\
\env{LESSKEYIN}, \env{LESSKEY\_CONTENT} & Key bindings. \\
\env{FASTLESSHISTFILE} & The history file (\file{\tl/.fastlesshst}; \cmd{-} for none). \\
\env{TMPDIR} & Temporary files, without \opt{--tmpdir}. \\
\end{reftable}

\subsection{Exit status}

\begin{reftable}[0.12\linewidth]
\thead{Status}{Meaning}
0 & Success; for \opt{--find}, something matched. \\
1 & \opt{--find} found nothing. \\
2 & An error, such as an unreadable file or damaged compressed data, or \kbd{\caret C} with \opt{-K}. \\
130 & Interrupted by \kbd{\caret C} outside the pager, as while decoding. \\
128+$n$ & Ended by signal $n$, such as SIGTERM. \\
\end{reftable}

% ---------------------------------------------------------------- differences
\section{Differences from GNU less}\label{sec:less}

\fl{} follows \cmd{less} 668 wherever it can, so muscle memory, \env{LESS}
settings, and lesskey files carry over. These are the places where it does not.

\subsection{What fastless adds}

\begin{reftable}[0.27\linewidth]
\thead{Feature}{Details}
Search modes & Literal and whole-line glob searches beside regular expressions
  (\opt{--literal}, \opt{--glob}, \kbd{ESC-r} \kbd{ESC-l} \kbd{ESC-w}). \\
Built-in decoding & gzip, zstd, xz, and bzip2, detected by content, decoded on several
  threads, verified; \opt{--compression}, \opt{--max-mem}, \opt{--tmpdir}. \\
Parallel search & Long searches and line counts use every core
  (\opt{--threads}); literal matching uses AVX2 (\opt{--simd}). \\
Byte positions & \kbd{P} jumps to a byte offset, and \opt{--offset} starts there. \\
Batch search & \opt{--find} and \opt{--all} print match offsets without the pager. \\
Status line & A default status line with file, position, and search mode. \\
\end{reftable}

\subsection{What behaves differently}

\begin{reftable}[0.27\linewidth]
\thead{Area}{In fastless}
Long lines & Always chopped, never wrapped: \opt{-S} is the only behaviour, and
  \opt{--wordwrap} is not supported. \\
Line numbers & Counted lazily, only when something shows them; a long jump does not
  wait for the count. \cmd{less} counts as it reads. \\
Regular expressions & POSIX extended expressions from the C library in the C
  locale, so case folding is ASCII only. Literal searches also match across line
  ends. \\
Positions & Offsets, percentages, and the status line count bytes of the decoded
  data. \\
Preprocessors & Compressed files bypass \env{LESSOPEN}. Standard input is never
  preprocessed, even with the \cmd{|-} form. \\
Control characters & \opt{-r} is the same as \opt{-R}: only colour sequences reach the
  terminal, and \opt{-u} never sends raw backspaces or carriage returns. Other bytes
  show as \cmd{\caret X}, \cmd{<XX>}, or \cmd{<U+XXXX>}; \env{LESSBINFMT} is not used. \\
\opt{-?} & Prints an option summary and exits, rather than opening the help screen. \\
History & Kept in \file{\tl/.fastlesshst} (\env{FASTLESSHISTFILE}), apart from
  \cmd{less}'s. \\
\end{reftable}

\subsection{What is not supported}

\begin{itemize}
\item Tags (\opt{-t}, \opt{-T}, \cmd{:t}), log files (\opt{-o}, \opt{-O}), colours
  (\opt{-D}, \opt{--use-color}), \opt{--modelines}, \opt{-"},
  \opt{--use-backslash}, OSC~8 hyperlinks, and compiled lesskey files (\opt{-k}).
\item The \opt{-}\meta{n} form of \opt{-z}\meta{n} (as in \opt{-20}), and
  capitalising only the first letter of an upper-case long option (as in
  \opt{--Quit-at-eof}).
\item Environment variables not listed above, such as \env{LESSHISTFILE},
  \env{LESSANSIENDCHARS}, \env{LESSMETACHARS}, \env{LESS\_IS\_MORE}, and the
  \env{LESS\_TERMCAP\_*} variables.
\item Character sets other than UTF-8 and single bytes.
\end{itemize}

% ---------------------------------------------------------------- performance
\section{Performance and limits}

\subsection{What to expect}

Opening a plain file costs nothing, whatever its size. Searching reads the bytes
it covers from storage, so a cold search runs at disk speed; once the file is
cached, literal searches run at several GiB per second. On an 8-core Ryzen 7
5700X with a warm cache, a case-insensitive literal search measured about
8.2\,GiB/s, and a regex without a usable literal ran 6.7 times faster on eight
threads than on one. Numbering the end of a 4\,GiB, 84-million-line file took
0.15\,s.

bzip2 decoding, 556\,MiB of logs compressed by \cmd{bzip2 -9}, wall time:

\begin{center}\small
\begin{tabular}{@{}rrrr@{}}
\toprule
\textsf{\bfseries Threads} & \textsf{\bfseries bzip2 (libbz2)} & \textsf{\bfseries lbzip2} & \textsf{\bfseries fastless} \\
\midrule
1 & 11.6\,s & 9.5\,s & 9.4\,s \\
4 & --- & 2.4\,s & 2.4\,s \\
16 & --- & 1.8\,s & 1.8\,s \\
\bottomrule
\end{tabular}
\end{center}

\subsection{Limits}

\begin{itemize}
\item A mapped file must not be truncated while it is shown, or the program may
  be killed by SIGBUS, as with any program that maps files. Replacing
  the file by name changes nothing unless \kbd{F} is used with
  \opt{--follow-name}.
\item A regular expression cannot match within a single line longer than
  2\,GiB; literal and glob searches can.
\item Archive members, such as those of a tar file, are not extracted; only the
  compression is removed.
\item \kbd{\caret C} is checked between pieces of work, so one very long line or
  one costly regular-expression call can delay it.
\item \fl{} needs a 64-bit system with POSIX \cmd{REG\_STARTEND}, such as Linux or
  macOS.
\end{itemize}

\subsection{Licence}

\fl{} is free software under the GNU General Public License, version 3. It includes
the bzip2 block decoder of lbzip2, which is licensed under the GPL, version 3 or
later.

\end{document}
